\section*{Chapter \ref{ch:g11:proving-algorithms-correct} --- Proving Algorithms Correct}

\begin{solution}{exo:g11:proving-algorithms-correct:1}
Turn $0$ swaps the smallest, $1$, into place $0$: \code{[1, 2, 9, 5, 7]}. Turn $1$ finds $2$ already in place $1$:
the list stays \code{[1, 2, 9, 5, 7]}.
\end{solution}

\begin{solution}{exo:g11:proving-algorithms-correct:2}
\code{n} itself, if \code{n} starts at $0$ or more: it falls by $2$ at each turn, and the loop runs only while
\code{n > 0}, so during the loop it is a whole number at least $1$. (If \code{n} starts odd, it ends at $-1$; the
\omterm{def:g11:proving-algorithms-correct:variant}{variant} only has to be non-negative while the loop runs.)
\end{solution}

\begin{solution}{exo:g11:proving-algorithms-correct:3}
\code{lo=0 mid=5 hi=10}: $30 < 52$; \code{lo=6 mid=8 hi=10}: $47 < 52$; \code{lo=9 mid=9 hi=10}: found, the answer is
$9$.
\end{solution}

\begin{solution}{exo:g11:proving-algorithms-correct:4}
Yes: before the first turn \code{i} is $0$ in the statement, \code{t[:0]} is empty, its sum is $0$, and \code{s} is
$0$.
\end{solution}

\begin{solution}{exo:g11:proving-algorithms-correct:5}
The end step. The invariant still holds at each turn, but the loop stops at \code{i = len(t) - 1}, so it only gives
``\code{best} is the largest of \code{t[:len(t) - 1]}'': the last element is never looked at, and $11$ is missed.
\end{solution}

\begin{solution}{exo:g11:proving-algorithms-correct:6}
\code{n} is a \omterm{def:g11:proving-algorithms-correct:variant}{variant}. Inside the loop \code{n > 1}; then \code{n // 2} and \code{n - 1} are both whole numbers smaller
than \code{n} and at least $1$. So \code{n} stays non-negative and strictly decreases: the loop stops.
\end{solution}

\begin{solution}{exo:g11:proving-algorithms-correct:7}
\smallskip\noindent\hspace*{2em}\code{def index_of_min(t, i):}

\noindent\hspace*{4em}\code{m = i}

\noindent\hspace*{4em}\code{for j in range(i + 1, len(t)):}

\noindent\hspace*{6em}\code{if t[j] < t[m]:}

\noindent\hspace*{8em}\code{m = j}

\noindent\hspace*{4em}\code{assert i <= m < len(t) and t[m] == min(t[i:])}

\noindent\hspace*{4em}\code{return m}

\smallskip\noindent The invariant of the loop is ``\code{t[m]} is a smallest element of \code{t[i:j]}''.
\end{solution}

\begin{solution}{exo:g11:proving-algorithms-correct:8}
The keep step of the invariant. When \code{t[mid] > x}, every element of \code{t[mid:]} is larger than \code{x}, but
\code{t[mid - 1]} may be \code{x}; \code{hi = mid - 1} puts it outside the window, so ``\code{x} is not in
\code{t[hi:]}'' becomes false. The \omterm{def:g11:proving-algorithms-correct:variant}{variant} still falls, which is why the loop stops, with a wrong answer.
\end{solution}

\begin{solution}{exo:g11:proving-algorithms-correct:9}
Invariant: at the start of the turn for \code{i}, \code{s} is the sum of \code{t[:i]}. Start: exercise
\ref{exo:g11:proving-algorithms-correct:4}. Keep: the turn adds \code{t[i]}, so \code{s} becomes the sum of
\code{t[:i + 1]}, which is the invariant for the next \code{i}. End: after the last turn, \code{i} has gone through
\code{len(t) - 1}, so \code{s} is the sum of \code{t[:len(t)]}, the whole list. A \code{for} loop over a range stops.
\end{solution}

\begin{solution}{exo:g11:proving-algorithms-correct:10}
\smallskip\noindent\hspace*{2em}\code{c = 0}

\noindent\hspace*{2em}\code{for i in range(len(t)):}

\noindent\hspace*{4em}\code{if t[i] == x:}

\noindent\hspace*{6em}\code{c = c + 1}

\smallskip\noindent Invariant: at the start of the turn for \code{i}, \code{c} is the number of occurrences of
\code{x} in \code{t[:i]}. The turn adds $1$ exactly when \code{t[i]} is \code{x}, so it is kept; at the end \code{c}
counts the whole list.
\end{solution}

\begin{solution}{exo:g11:proving-algorithms-correct:11}
\smallskip\noindent\hspace*{2em}\code{def is_sorted(t):}

\noindent\hspace*{4em}\code{for i in range(1, len(t)):}

\noindent\hspace*{6em}\code{if t[i - 1] > t[i]:}

\noindent\hspace*{8em}\code{return False}

\noindent\hspace*{4em}\code{return True}

\smallskip\noindent Invariant: at the start of the turn for \code{i}, \code{t[:i]} is sorted. It holds for
\code{i = 1}, a single element. If \code{t[i - 1] > t[i]}, the list is not sorted and \code{False} is right;
otherwise \code{t[i]} is not smaller than the last element of the sorted \code{t[:i]}, so \code{t[:i + 1]} is sorted.
If the loop ends without returning, \code{t[:len(t)]} is sorted and \code{True} is right. An empty list or a list of
one element skips the loop and is sorted.
\end{solution}

\begin{solution}{exo:g11:proving-algorithms-correct:12}
Outer invariant: at the start of turn \code{i}, \code{t[:i]} holds the first \code{i} elements of the original list,
sorted, and \code{t[i:]} is unchanged. Start: \code{t[:1]} has one element. Inner invariant, with $x$ the old
\code{t[i]}: \code{t[:j]} is unchanged and sorted, \code{t[j + 1:i + 1]} holds the old \code{t[j:i]}, every element of
which is larger than $x$. It holds when \code{j = i}, both slices being as they were. An inner turn copies
\code{t[j - 1]}, larger than $x$, into place \code{j} and lowers \code{j}: kept. The inner loop stops, \code{j} being a
\omterm{def:g11:proving-algorithms-correct:variant}{variant} that is at least $0$. When it stops, either \code{j = 0}, or \code{t[j - 1] <= x}; writing $x$ in place
\code{j} puts it after elements not larger than it and before elements larger than it: \code{t[:i + 1]} is sorted and
holds the old \code{t[:i]} and $x$. The outer invariant is kept; at the end, \code{i} reaches \code{len(t)} and the whole
list is sorted. The outer loop is a \code{for} loop and stops.
\end{solution}

\begin{solution}{exo:g11:proving-algorithms-correct:13}
Let $p$ be an index with \code{t[p] == x}. By the invariant, \code{x} is not in \code{t[:lo]} nor in \code{t[hi:]}, so
$p$ always lies in the window: \code{lo <= p < hi}, and the window is never empty. The loop condition \code{lo < hi}
therefore stays true, and the only way out is the \code{return mid}, which happens only when \code{t[mid] == x}. Since
the \omterm{def:g11:proving-algorithms-correct:variant}{variant} \code{hi - lo} forces the loop to stop, the function does return, and what it returns is an index of
\code{x}.
\end{solution}

\begin{solution}{exo:g11:proving-algorithms-correct:14}
\code{n} does not always fall: an odd \code{n} becomes \code{3 * n + 1}, much larger. ``The number of turns left''
falls by one each turn, but it is a whole number only if the loop stops, which is what was to be proved: using it
assumes the result. A proof would need some quantity, computed from \code{n}, that never becomes negative and falls at
each turn, or a different argument altogether. Nobody has found one: whether this loop stops for every \code{n} is an
open question, though every number tried so far reaches $1$.
% ledger: g11:collatz-open
\end{solution}

\begin{solution}{exo:g11:proving-algorithms-correct:15}
\smallskip\noindent\hspace*{2em}\code{def first_true(f, lo, hi):}

\noindent\hspace*{4em}\code{while lo < hi:}

\noindent\hspace*{6em}\code{mid = (lo + hi) // 2}

\noindent\hspace*{6em}\code{if f(mid):}

\noindent\hspace*{8em}\code{hi = mid}

\noindent\hspace*{6em}\code{else:}

\noindent\hspace*{8em}\code{lo = mid + 1}

\noindent\hspace*{4em}\code{return lo}

\smallskip\noindent Invariant: \code{f} is false on every number from the first \code{lo} to the present
\code{lo} excluded, and true from the present \code{hi} to the first \code{hi} excluded. A true \code{f(mid)} makes
\code{f} true from \code{mid} on, by the \omterm{def:g10:specifications-and-tests:pre-post}{precondition}; a false one makes it false up to \code{mid}. At the end
\code{lo == hi}: everything before is false, everything from there is true, so \code{lo} is the first true, or the
first \code{hi} if none is. \omterm{def:g11:proving-algorithms-correct:variant}{Variant}: \code{hi - lo}, which both branches lower since \code{lo <= mid < hi}.
\end{solution}

\begin{solution}{pb:g11:proving-algorithms-correct:1}
\textbf{1.} \omterm{def:g10:specifications-and-tests:pre-post}{Precondition}: $n$ is a whole number, $n \ge 0$. \omterm{def:g10:specifications-and-tests:pre-post}{Postcondition}: $r$ is a whole number, $r^2 \le n$, and no
larger whole number has its square at most $n$.

\textbf{2.} For $r \ge 0$, the squares grow with $r$. So if $(r + 1)^2 > n$, every number above $r$ has a square
above $n$, and $r$ is the largest; conversely the largest $r$ has $(r + 1)^2 > n$.

\textbf{3.} $0$, $1$, $9$, $10$, $10$.

\textbf{4.} $r = 0$ has $r^2 \le n$, so there is at least one candidate; a candidate $r \ge 1$ has $r \le r^2 \le n$,
so there are finitely many. Their largest exists and is unique.

\textbf{5.} $r^2 \le n$. When the loop stops, its condition is false: $(r + 1)^2 > n$, and with question 2 that is the
\omterm{def:g10:specifications-and-tests:pre-post}{postcondition}.

\textbf{6.} Start: $r = 0$ and $0 \le n$. Keep: a turn runs only if $(r + 1)^2 \le n$, and then sets $r$ to $r + 1$,
whose square is at most $n$.

\textbf{7.} By the invariant, \code{n - r * r} is at least $0$; each turn raises $r$ by one, so $r^2$ grows by
$2r + 1 \ge 1$ and \code{n - r * r} falls by at least $1$.

\textbf{8.} One turn per value of $r$ above $0$: as many turns as the result, $1000$ for $n = 10^6$, as printed.

\textbf{9.} \code{lo = 0} and $0 \le n$; \code{hi = n + 1} and $(n + 1)^2 \ge n + 1 > n$.

\textbf{10.} If \code{mid * mid <= n}, then \code{lo = mid} has its square at most $n$, and \code{hi} does not change.
Otherwise $n < \mathit{mid}^2$, and \code{hi = mid} keeps $n < \mathit{hi}^2$.

\textbf{11.} At the end \code{hi - lo <= 1}. By the invariant $\mathit{lo}^2 \le n < \mathit{hi}^2$, so \code{lo < hi},
hence \code{hi = lo + 1} and $\mathit{lo}^2 \le n < (\mathit{lo} + 1)^2$: the \omterm{def:g10:specifications-and-tests:pre-post}{postcondition}, by question 2.

\textbf{12.} Inside the loop \code{hi - lo >= 2}, so $\mathit{lo} + 1 \le (\mathit{lo} + \mathit{hi}) / 2 \le
\mathit{hi} - 1$, and rounding down keeps \code{mid} between \code{lo + 1} and \code{hi - 1}.

\textbf{13.} \code{hi - lo} is a whole number, at least $2$ inside the loop; each turn replaces \code{lo} or
\code{hi} by \code{mid}, strictly between them, so it strictly falls. The loop stops.

\textbf{14.} \code{lo=0, hi=16}: \code{mid=8}, $64 > 15$, \code{hi=8}; \code{mid=4}, $16 > 15$, \code{hi=4};
\code{mid=2}, $4 \le 15$, \code{lo=2}; \code{mid=3}, $9 \le 15$, \code{lo=3}; now \code{hi - lo = 1}. Four turns,
result $3$, as printed.

\textbf{15.} The new window is \code{mid - lo} or \code{hi - mid}; these two add up to \code{hi - lo} and differ by at
most $1$, so the larger is at most half of \code{hi - lo} rounded up. Half of $2^k$ is $2^{k - 1}$, already whole.

\textbf{16.} The window starts at $n + 1 \le 2^k$. After one turn it is at most $2^{k - 1}$, after two at most
$2^{k - 2}$, and after $k$ turns at most $2^0 = 1$, where the loop stops: at most $k$ turns.

\textbf{17.} $2^{20} = 1\,048\,576 \ge 10^6 + 1$: at most $20$ turns for $n \le 10^6$.

\textbf{18.} The run finds $20$ turns at most, first at $n = 555\,007$. For $n + 1 \le 2^{19}$ the bound gives $19$;
just above, $20$ turns are possible but not forced: whether the window shrinks to the half rounded down or rounded up
depends on the branch taken, and only from $555\,007$ do enough halvings round up.

\textbf{19.} The assertion compares a million results with Python's own function, an independent computation: it would
catch a slip between the code and the proof, such as a wrong sign. The proof covers every $n$, beyond $10^6$ too, and
\omterm{def:g11:proving-algorithms-correct:variant}{termination}, which no test run can show.

\textbf{20.} The integer square root by halving is \textbf{\omterm{def:g11:proving-algorithms-correct:correct}{correct} for every whole $n$ and always stops}; for
$n \le 10^6$ it takes \textbf{at most $20$ turns}, reached first at $n = 555\,007$, where counting up takes up to
$1000$.
\end{solution}
